\section{Pipeline details}
\label{app:pipeline}

Lab2Sim is implemented in Scenario Forge as a producer-side agent in front of
the existing package compiler. Request understanding, search, specification,
and generation are the new front-end. ConvertAsset, source bindings, asset
locks, and scenario compile are unchanged.

\paragraph{Request.}
The agent reads the experiment requirement and emits an asset request list:
one row per object, with a slug, why the experiment needs it, whether it
articulates, which dimensions are task-critical, and the intended robot
contact (grasp, pour, open, place). Text, stills, and video are parsed into
that same table.

\paragraph{Evidence.}
For each row the agent retrieves manufacturer identity. Catalog pages,
datasheets, and official photographs are hashed into an evidence register.
Items are graded as exact-model evidence, family evidence, derived crops, or
unknown. Unknown fields stay unknown; they are not filled in from pixels.

\paragraph{Specification.}
The specification template has a manufacturer table, a decomposition
decision, a prim tree, joint contracts, state variables, interaction
semantics, a QA matrix, delivery conventions, and an open-measurement punch
list. Table~\ref{tab:oven-spec} is the manufacturer table of the OVEN~125
specification behind Figure~\ref{fig:real-spec}a. The burette specification
fixes a revolute stopcock on \texttt{burette.usd}, a static
\texttt{burette\_stand.usd}, and a scene-level fixed joint between their mount
frames.

\paragraph{Generation.}
The skill iterates geometry and articulation against the specification.
Blender is the always-on modeller. Hunyuan3D is called only when the service
is up; a miss falls back to Blender without changing the specification. A
static articulation checker runs under Isaac Kit Python. Failures block
admission.

\paragraph{Admission.}
ConvertAsset snapshots the generated tree, converts meshes and materials,
and qualifies PhysX. A promoted asset carries a promotion receipt with its
hash and the qualifying runtime; the scene that uses it carries a render
manifest with camera poses and frame hashes. Figure~\ref{fig:real-spec} is
built from those two records for each row: the OVEN~125 asset was promoted
under Isaac Sim~4.1 with a 4.5 compatibility check, and the burette station
under Isaac Sim~4.5. Lab2Sim does not own a second physics stack.

\section{OVEN~125 specification excerpt}
\label{app:spec}

\begin{table}[h]
  \centering
  \small
  \caption{Manufacturer fields for the IKA OVEN~125 control-dry (solid door)
  as fixed in the agent-written asset specification.}
  \label{tab:oven-spec}
  \begin{tabular}{@{}llp{3.4cm}@{}}
    \toprule
    Field & Value & Role \\
    \midrule
    Envelope $W{\times}H{\times}D$ & $700{\times}825{\times}650$~mm & Body collider \\
    Chamber $W{\times}D{\times}H$ & $550{\times}525{\times}450$~mm & Interior volume \\
    Volume & 125~L & Identity \\
    Mass & 82~kg & Fixed base \\
    Door travel & ${\sim}180^{\circ}$ revolute & Open / close \\
    Hinge & Configurable side & Joint axis \\
    Shelves & Up to 6; 30~kg each & Placement surfaces \\
    Controls & Dual TFT, rotary/push knob & Setpoint \\
    \bottomrule
  \end{tabular}
\end{table}

Decomposition follows the specification rather than the mesh.
\texttt{oven\_body} is a fixed base. \texttt{door} is a revolute link.
\texttt{handle} is a grasp collider on the door.
\texttt{shelf\_0}--\texttt{shelf\_5} carry placement surfaces.
\texttt{control\_knob} is a rotary/push compound. Generation may rebuild
geometry. It may not rename those contracts.

\section{Additional renders}
\label{app:extra}

Figure~\ref{fig:app-extra} adds three close views. The oven interior shows
the lower shelf of the admitted OVEN~125 (r2), with a beaker and flask seated
on the wire deck. The powder close-up is the weighing scoop of the r4 station
in Figure~\ref{fig:hero}. The Lift2 frame is the same Isaac runtime that
consumes the compiled laboratories: a bimanual robot already spawned in front
of a table.

\begin{figure*}[h]
  \centering
  \begin{tabular}{@{}ccc@{}}
    \includegraphics[width=0.31\textwidth]{source/app_oven_interior.png} &
    \includegraphics[width=0.31\textwidth]{source/powder_r4_hero.png} &
    \includegraphics[width=0.31\textwidth]{source/app_lift2_runtime.png} \\
    {\small Oven interior} &
    {\small Powder scoop} &
    {\small Lift2 in Isaac} \\
  \end{tabular}
  \caption{Extra Isaac frames. Left: OVEN~125 lower shelf with glassware.
  Center: powder-weighing scoop. Right: Lift2 in the Isaac runtime that
  loads Lab2Sim scenes.}
  \label{fig:app-extra}
\end{figure*}

\section{Figure sources}
\label{app:figures}

Figures~\ref{fig:hero}--\ref{fig:real-spec} use frames from the
current-delivery scenes and from the indexed manufacturer evidence. Figure~2
shows eBench task05 (r11), task07 (r10.1), the stir-bar station (r5), and
the water-bath station (r1). Figure~3 uses the OVEN~125 Task~12 scene (r2)
and the titration station (r1.7); the open-oven frame is a camera-only
retake of the admitted scene chosen to match the catalog viewpoint. Frames
are cropped and, where noted in the figure provenance, globally
brightness-adjusted; scene state, geometry, materials, and physics are never
edited for a figure. Every scene in the paper can be generated with the
Blender skill alone.
